% 图 18.1 —— 由 `checks/emit_hand.py` 自动拼出（probe 精测 + prims2 图元分解）
%%
% ⚠ 这是**稿子不是成品**：跑 `checks/checkhand.py 18.1` 看指标 + 看叠合图。
%%
%% ⚠ 待核：
%   · 本图 probe 报出阴影族 1 个 —— **本脚本不生成阴影**（要照原书手写）；prims2 的 strokes 里可能含阴影线，核对时留意别画重
%   · 可疑字形 1 项（空心圈 ⊙ 常被认成字母 o）—— 见文件末
%%
\documentclass[12pt]{standalone}
\usepackage{fontspec}
\setsansfont{Arial}
\newfontfamily\mfallback{DejaVu Sans}
\usepackage{tikz}
\begin{document}
\begin{tikzpicture}[x=1pt, y=-1pt, line width=6.0pt, line cap=round, line join=round]

  %% ════ 填充区（prims2 的 fills）════

  %% ════ 直线段（prims2 的 strokes）════
  \draw[line width=6.0pt] (628.0,79.0) -- (612.0,79.0);
  \draw[line width=5.0pt] (236.0,216.0) -- (236.0,224.0);
  \draw[line width=6.0pt] (607.0,85.0) -- (610.0,80.0);
  \draw[line width=4.0pt] (609.0,79.0) -- (465.0,79.0);
  \draw[line width=7.7pt] (475.0,134.0) -- (470.0,130.0);
  \draw[line width=5.0pt] (945.0,162.0) -- (964.9,147.1);
  \draw[line width=5.0pt] (964.9,147.1) -- (986.0,133.0);
  \draw[line width=5.0pt] (966.0,22.0) -- (931.3,21.0);
  \draw[line width=4.0pt] (931.3,21.0) -- (904.1,28.0);
  \draw[line width=4.0pt] (690.7,24.3) -- (684.0,33.5);
  \draw[line width=5.0pt] (684.0,33.5) -- (684.0,42.5);
  \draw[line width=5.0pt] (371.0,125.0) -- (369.0,206.9);
  \draw[line width=4.0pt] (369.0,206.9) -- (360.6,219.4);
  \draw[line width=5.0pt] (235.0,225.0) -- (132.2,228.0);
  \draw[line width=4.0pt] (46.0,58.8) -- (47.5,51.5);
  \draw[line width=4.0pt] (47.5,51.5) -- (83.6,42.0);
  \draw[line width=5.0pt] (83.6,42.0) -- (109.8,40.0);
  \draw[line width=4.0pt] (615.0,129.0) -- (471.0,129.0);
  \draw[line width=5.0pt] (469.0,128.0) -- (452.0,129.0);

  %% ════ 曲线（prims2 的 curves —— 三段式，坐标同为 (y,x)）════
  \draw[line width=4.0pt] (370.0,124.0) .. controls (314.6,127.2) and (253.7,160.1) .. (236.0,216.0);
  \draw[line width=4.0pt] (801.0,211.0) .. controls (838.3,221.9) and (860.8,274.1) .. (904.2,262.8);
  \draw[line width=5.0pt] (904.2,262.8) .. controls (942.4,241.8) and (920.7,192.2) .. (944.0,163.0);
  \draw[line width=4.0pt] (986.0,133.0) .. controls (1032.8,105.5) and (1016.1,32.9) .. (966.0,22.0);
  \draw[line width=5.0pt] (904.1,28.0) .. controls (866.5,45.4) and (821.5,34.9) .. (782.4,28.0);
  \draw[line width=5.0pt] (782.4,28.0) .. controls (755.7,13.9) and (718.7,10.1) .. (690.7,24.3);
  \draw[line width=4.0pt] (684.0,42.5) .. controls (700.5,89.4) and (661.6,136.0) .. (681.0,182.0);
  \draw[line width=4.0pt] (681.0,182.0) .. controls (688.8,200.4) and (706.8,212.2) .. (725.5,215.0);
  \draw[line width=4.0pt] (725.5,215.0) .. controls (750.8,215.1) and (774.3,207.1) .. (800.0,211.0);
  \draw[line width=5.0pt] (360.6,219.4) .. controls (323.1,233.8) and (277.7,223.5) .. (237.0,225.0);
  \draw[line width=5.0pt] (132.2,228.0) .. controls (107.8,236.5) and (76.3,246.6) .. (52.6,229.6);
  \draw[line width=4.0pt] (52.6,229.6) .. controls (35.1,214.4) and (33.4,191.6) .. (33.0,169.7);
  \draw[line width=4.0pt] (33.0,169.7) .. controls (41.6,134.2) and (62.9,95.6) .. (46.0,58.8);
  \draw[line width=5.0pt] (109.8,40.0) .. controls (201.4,47.3) and (289.3,12.3) .. (376.9,37.0);
  \draw[line width=4.0pt] (376.9,37.0) .. controls (386.6,42.9) and (398.1,49.3) .. (402.0,60.6);
  \draw[line width=5.0pt] (402.0,60.6) .. controls (406.5,86.5) and (377.6,101.3) .. (371.0,124.0);
  \draw[line width=4.0pt] (944.0,162.0) .. controls (894.6,140.4) and (827.0,160.9) .. (801.0,210.0);

  %% ════ 实心点 ════
  \fill (611.0,74.0) circle (4.1);

  %% ════ 标签（probe 直接给的 \node 行）════
  %% ⚠ 全部补了 fill=white：文字在最上层，否则与曲线交叉干扰
     \node[anchor=center,inner sep=0pt,fill=white,font=\fontsize{27.8}{34.8}\selectfont] at (546.0,54.0) {\textsf{\textbf{\textit{f}}}};   % box(540,43,11,21)
     \node[anchor=center,inner sep=0pt,fill=white,font=\fontsize{37.2}{46.5}\selectfont] at (29.2,109.5) {\textsf{\textbf{\textit{x}}}};   % box(18,96,20,22)
     \node[anchor=center,inner sep=0pt,fill=white,font=\fontsize{30.3}{37.9}\selectfont] at (1035.7,119.0) {\textsf{\textbf{\textit{V}}}};   % box(1026,107,20,22)
     \node[anchor=center,inner sep=0pt,fill=white,font=\fontsize{23.7}{29.7}\selectfont] at (558.8,152.4) {\textsf{\textbf{1}}};   % box(552,144,8,16)
     \node[anchor=center,inner sep=0pt,fill=white,font=\fontsize{27.8}{34.8}\selectfont] at (529.0,160.0) {\textsf{\textbf{\textit{f}}}};   % box(523,149,11,21)
     \node[anchor=center,inner sep=0pt,fill=white,font=\fontsize{26.5}{33.2}\selectfont] at (550.9,156.0) {{\mfallback\textbf{−}}};   % box(538,150,13,4)
     \node[anchor=center,inner sep=0pt,fill=white,font=\fontsize{29.6}{37.0}\selectfont] at (317.0,186.2) {\textsf{\textbf{\textit{U}}}};   % box(307,174,20,22)
     \node[anchor=center,inner sep=0pt,fill=white,font=\fontsize{27.1}{33.9}\selectfont] at (866.5,199.5) {\textsf{\textbf{\textit{f}}}};   % box(861,188,10,22)
     \node[anchor=center,inner sep=0pt,fill=white,font=\fontsize{30.7}{38.4}\selectfont] at (878.5,202.5) {\textsf{\textbf{(}}};   % box(874,188,8,28)
     \node[anchor=center,inner sep=0pt,fill=white,font=\fontsize{29.6}{37.0}\selectfont] at (895.0,200.2) {\textsf{\textbf{\textit{U}}}};   % box(885,188,20,22)
     \node[anchor=center,inner sep=0pt,fill=white,font=\fontsize{30.7}{38.4}\selectfont] at (914.5,202.5) {\textsf{\textbf{)}}};   % box(909,188,8,28)

  % 钉死画布：否则超出的图元/控制点会把页面撑大，1:1 回比就失准
  \pgfresetboundingbox
  \path[use as bounding box] (0,0) rectangle (1055,275);
\end{tikzpicture}
\end{document}

% ⚠ 待核清单（probe 拿不准，人眼过一遍）：
%   · 未识别/跳过：⑤ 跳过/未识别的字形
